% SOURCE AUTHORITY:
% expose_VIII_body.tex lines 1976--2003; scan indices 298--300;
% source folios 287--289.
% Direct scan comparison: physical PDF pages 299--301 in overlapping
% 1100-dpi bands.  The prose, formulas, and subscripts are fully legible at
% that resolution; no higher-resolution escalation was needed in this unit.
% Transparent source dispositions:
% - the stray \mathbb Z_Z in the paragraph preceding Lemma 6.2 is corrected
%   to \mathbb Z_S, the free group scheme over the stated base S;
% - the corrupt insertion between "over S" and \overline G in Lemma 6.2 is
%   omitted; the authority clearly defines alpha:S->\mathbb Z_{S,T} and then
%   pulls back the extension \overline G along alpha.
% Hard stop before 6.2.1, source line 2004 / scan index 300.

\medskip
\noindent\underline{6.  Extensions and biextensions by the Neron model of
$G_m$.}

\medskip
\noindent 6.1.  Let $S$ be a scheme, $T$ a closed subscheme, and $Z$ a set,
giving a truncated constant scheme $Z_{S,T}$ over $S$ (5.7).  There is a
canonical homomorphism, functorial in $Z$, of etale schemes over $S$,
\begin{equation}\tag{6.1.1}
 Z_S\longrightarrow Z_{S,T},
\end{equation}
which is a \underline{surjective} morphism, as follows at once from (5.7.1).
Its functoriality in $Z$ shows that when $Z$ is a group, (6.1.1) is a
homomorphism of etale group schemes over $S$.  Its kernel is therefore an
etale group scheme over $S$, denoted $Z_{S,U}$, and there is a canonical
exact sequence, functorial in $Z$,
\begin{equation}\tag{6.1.2}
 0\longrightarrow Z_{S,U}\longrightarrow Z_S\longrightarrow Z_{S,T}
 \longrightarrow 0.
\end{equation}
It follows immediately from (5.7.1) that $Z_{S,U}$ is the sum of a family of
$S$-schemes $S_n$, $n\in Z$, where $S_n=S$ for $n=0$ and $S_n=U$ for
$n\ne0$.  This description of $Z_{S,U}$ readily shows that for every group
$F$ over $S$, restriction to $U$ induces a bijection
\begin{equation}\tag{6.1.3}
 \mathrm{Hom}_{\mathrm{gr}}(Z_{S,U},F)
 \xrightarrow{\ \sim\ }
 \mathrm{Hom}_{\mathrm{gr}}(Z_U,F|U).
\end{equation}

\medskip
\noindent 6.1.4.  Henceforth we suppose that $Z$ is commutative, and we wish
to examine extensions of commutative groups of $Z_{S,T}$ by a given
commutative group $G$ over $S$.  By the exact sequence (6.1.2), the category
of these extensions is equivalent to the category of extensions $F$ of
$Z_S$ by $G$, furnished with a trivialization of the inverse image of the
extension along $Z_{S,U}\to Z_S$.  By (6.1.3), such a trivialization amounts
to a trivialization of the extension $F|U$ of $Z_U$ by $G|U$.

For what follows we take $Z$ to be the group $\mathbb Z$ of integers.  Thus
$\mathbb Z_S$ is the free group on one generator, and VII 1.4 gives an
equivalence between the category of extensions of $\mathbb Z_S$ by $G$ and
the category of torsors under $G$.  Applying the same result to
$\mathbb Z_U$ and $G|U$ gives the following.

\medskip
\noindent\underline{Lemma 6.2.}  \underline{Let $S$ be a scheme, $T$ a
closed subscheme, $U=S\setminus T$, and $G$ a commutative group over $S$.
The category of commutative groups over $S$ which are extensions of
$\mathbb Z_{S,T}$ (defined in (5.7)) by $G$ is equivalent to the category of
pairs $(F,s)$ consisting of a torsor $F$ under $G$ and a section $s$ of
$F|U$.  If $\alpha:S\to\mathbb Z_{S,T}$ is the section of
$\mathbb Z_{S,T}$ which is the image of the section $1$ of $\mathbb Z_S$
over $S$, then the torsor corresponding to an extension $\overline G$ of
$\mathbb Z_{S,T}$ by $G$ is $F=\alpha^*(\overline G)$ under $G$, furnished
with the section of $F|U$ given by restricting the unit section of
$\overline G$ to $U$.}

